\section{Background and Related Work}\label{sec:background}

Executable architecture enforcement tools differ in how they encode constraints, but all aim to operationalize high-level architectural intent as checkable rules. ArchUnit, developed for Java and the JVM, allows assertions over classes, packages, and annotations through Java-based test specifications. Rule declarations generally capture dependency direction, package layering, inheritance restrictions, and annotation-based conventions. import-linter provides a Python-specific mechanism for checking dependencies between modules or packages, with built-in support for layering and import restrictions. dependency-cruiser offers similar functionality for JavaScript and TypeScript projects, identifying cycles, forbidden dependencies, and path-based exceptions. In each ecosystem, the mechanism is similar: architecture is encoded as executable rules that are executed in CI, often with exemptions and ignore lists to support exceptional cases.

The native exemption mechanisms are central to understanding rule decay. In ArchUnit, rules may be frozen with commands such as ``FreezingArchRule'' or ignored with annotations such as ``@ArchIgnore''. In import-linter, project maintainers can define import boundaries and exceptions, or use rules such as ``ignore\_imports'' to suppress specific paths. dependency-cruiser supports ignore patterns, exceptions, path filters, and path-based exceptions such as ``pathNot'' to permit targeted deviations. These escape hatches are practical and often necessary, but they also create an empirical risk: a rule set can look active while quietly becoming permissive, undermining the architectural guarantees it was designed to enforce.

Related work on architectural conformance and evolution has grown in recent years, yet it remains fragmented. Pontillo et al. study the role of architecture-as-code in practice and highlight tooling, documentation, and governance gaps~\cite{pontillo2026icsa}. Lucena et al. provide a taxonomy of design rules and show how rule intent can vary by architecture style and technology stack~\cite{lucena2026ecsa}. Larsen and Moghaddam focus on the longitudinal evolution of architectural governance in open source and document evidence of rule drift and adoption heterogeneity~\cite{larsen2026seaa}. Biazotto et al. analyze dependency management in modern JavaScript ecosystems and provide evidence that dependency boundaries and governance policies are often configured in ad hoc ways~\cite{biazotto2026}. Beyond these studies, work on architecture erosion emphasizes that architecture decay is not merely a function of code churn but also of governance and enforcement mechanisms~\cite{li2015,rinne2023,wagner2022}.

Table~\ref{tab:positioning} summarizes the position of our work relative to this literature. Existing studies have improved understanding of specific tools or specific ecosystems, but no prior work scales across multiple ecosystems and trajectories while explicitly operationalizing rule rot. Our study differs by mining thousands of repositories, reconstructing rule histories, and evaluating adoption effects with causal inference. We triangulate the rule-level and repository-level evidence with tool-native enforcement semantics, yielding a more complete picture of the gap between intended and realized architecture governance.

\begin{table}[t]
\centering
\caption{Positioning relative to prior architecture governance and AaC work.}
\label{tab:positioning}
\begin{tabular}{p{2.7cm}p{2.2cm}p{2.2cm}p{2.2cm}}
\toprule
\textbf{Study} & \textbf{Scope} & \textbf{Method} & \textbf{Contribution} \\
\midrule
Pontillo et al. & Industry AaC practice & Qualitative / multi-case & Adoption themes and gaps \\
Lucena et al. & NL design rules & Taxonomy + generation & Conceptual rule categories \\
Larsen and Moghaddam & OSS architectural quality & Longitudinal + causal & Smell trajectories under AI adoption \\
Biazotto et al. & CI/CD debt tooling & Pipeline mining & CI-based governance patterns \\
This paper & Multi-ecosystem OSS rules & Mining + history + cohort & Prevalence/CI map, \rrot, ArchRules \\
\bottomrule
\end{tabular}
\end{table}

